\subsubsection{Subproceso 1: Acreditación y Pre-requisitos}

\textbf{Personas:}
\begin{itemize}
    \item \textbf{Estudiante Regular:} Paga su autoseguro, regulariza sus deudas y verifica su habilitación y turno en el portal.
    \item \textbf{Personal de la DIRCE / ORCE:} Procesa el cierre de notas, las convalidaciones, el cálculo de promedios y la carga de turnos en el sistema.
    \item \textbf{Comisión de Matrícula, Dirección de Escuela, Tutores y Personal de Facultad:} Gestionan el calendario de turnos, cambios de plan, listas de promoción, pases de tutoría y reportes de deudas o pagos.
\end{itemize}

\textbf{Tecnología:}
\begin{itemize}
    \item \textbf{Portal Intralu (\url{https://intralu.uni.edu.pe/}):} Plataforma web donde el alumno genera y revisa sus órdenes de pago de autoseguro, consulta si registra deudas y verifica su avance curricular.
    \item \textbf{Intranet DIRCE (\url{https://dirce.uni.edu.pe/}):} Portal de matrícula donde se muestran la ficha de matrícula personalizada, los grupos y horarios de turno de la Facultad y la lista de requisitos validados (bloqueos, adeudos y pagos).
    \item \textbf{Sistema Integrado de Gestión Académica (SIGA-ORCE) y Sistema de Caja UNI:} Base de datos central que guarda las notas del ciclo anterior, calcula el orden de mérito y recibe la confirmación de pagos de recaudación.
\end{itemize}

\textbf{Documentos y referencias:}
\begin{itemize}
    \item \textbf{Texto Único de Procedimientos Administrativos (TUPA UNI)}~\cite{tupa2024}: Define los requisitos previos de la matrícula regular (pago de autoseguro y ausencia de adeudos en la Facultad).
    \item \textbf{Reglamento de Matrícula de Pregrado}~\cite{reglamento2022}: Regula el cierre previo de notas, el orden de mérito por promedio de los dos últimos ciclos, los cambios de plan y las condiciones de habilitación.
    \item \textbf{Reglamento del Fondo de Autoseguro Estudiantil}~\cite{autoseguro2023}: Establece el carácter obligatorio del aporte estudiantil en cada ciclo académico.
    \item \textbf{Ley Universitaria y Estatuto de la UNI}~\cite{ley30220}: Marco legal general de la condición de estudiante regular y del proceso de matrícula.
\end{itemize}
